\section{Evaluation Benchmarks and Metrics}
\label{sec:benchmarks}

Evaluating coding agents requires benchmarks that capture real-world utility. This section surveys major benchmarks and discusses their strengths and limitations.

\subsection{Function-Level Benchmarks}

\textbf{HumanEval}~\citep{chen2021codex} contains 164 hand-written Python programming problems. Each problem includes a function signature, docstring, and test cases. The model must generate a function body that passes all tests.

\textbf{HumanEval+}~\citep{liu2023humaneval} augments HumanEval with additional test cases to catch solutions that pass original tests through shortcuts or edge case failures.

\textbf{MBPP}~\citep{austin2021mbpp} (Mostly Basic Python Programming) contains 974 crowd-sourced Python problems designed to be solvable by entry-level programmers.

These benchmarks measure basic code generation ability but have significant limitations:
\begin{itemize}
    \item Problems are short, self-contained, and unrealistic
    \item No project context or cross-file dependencies
    \item Focus on algorithmic problems rather than software engineering
\end{itemize}

\subsection{Repository-Level Benchmarks}

\textbf{SWE-bench}~\citep{jimenez2024swebench} represents a paradigm shift in evaluation. It contains 2,294 real GitHub issues from popular Python repositories, paired with the pull requests that resolved them. An agent must:
\begin{enumerate}
    \item Understand the issue description
    \item Navigate an unfamiliar codebase
    \item Identify relevant files
    \item Generate a patch that resolves the issue
    \item Pass the repository's test suite
\end{enumerate}

Variants include:
\begin{itemize}
    \item \textbf{SWE-bench Lite}: 300-problem subset for faster evaluation
    \item \textbf{SWE-bench Verified}: 500 human-verified problems with confirmed solvability
\end{itemize}

SWE-bench has become the de facto standard for evaluating autonomous coding agents, though it has limitations:
\begin{itemize}
    \item Focuses primarily on bug fixes
    \item Python-only
    \item Test-based evaluation may miss code quality issues
\end{itemize}

\subsection{Other Benchmarks}

\textbf{CodeContests}~\citep{li2022codeforces} uses competitive programming problems, testing algorithmic reasoning rather than software engineering.

\textbf{DS-1000}~\citep{lai2022ds1000} focuses on data science tasks using pandas, numpy, and related libraries.

\textbf{CrossCodeEval}~\citep{ding2023crosscodeeval} evaluates cross-file context utilization.

\textbf{RepoBench}~\citep{liu2023repobench} measures repository-level code completion.

\subsection{Metrics}

\textbf{pass@k}: The probability that at least one of $k$ generated solutions passes all tests. Typically reported as pass@1, pass@10, pass@100.

\textbf{Resolution rate}: For SWE-bench, the percentage of issues where the generated patch passes the test suite.

\textbf{Cost per task}: Emerging metric tracking API costs, relevant for practical deployment.

\subsection{Current State of the Art}

\begin{table}[h]
\centering
\caption{Selected results on major benchmarks (as of early 2025)}
\label{tab:benchmarks}
\begin{tabular}{lcc}
\toprule
\textbf{Benchmark} & \textbf{Top Performance} & \textbf{System} \\
\midrule
HumanEval & 95\%+ pass@1 & GPT-4, Claude 3.5 \\
MBPP & 85\%+ pass@1 & Various frontier models \\
SWE-bench Verified & 40--55\% & Top agents (2024--2025) \\
\bottomrule
\end{tabular}
\end{table}

\subsection{Limitations of Current Evaluation}

Current benchmarks leave significant gaps:

\begin{itemize}
    \item \textbf{Data contamination}: Popular benchmarks may appear in training data
    \item \textbf{Narrow task types}: Most focus on bug fixes, missing feature development
    \item \textbf{Missing dimensions}: Code quality, maintainability, security not measured
    \item \textbf{Real-world gap}: Benchmark success may not predict production utility
\end{itemize}

\subsection{Emerging Approaches}

Researchers are exploring:
\begin{itemize}
    \item Human evaluation studies of AI-assisted development
    \item A/B testing in production environments
    \item Multi-dimensional quality assessment (correctness + style + security)
    \item Task diversity beyond bug fixing
\end{itemize}
